\section{Das Unternehmen und seine Organe}
\label{sec:unternehmen}

\subsection{Vorstand und Aufsichtsrat}

Der Vorstand besteht aus sieben Mitgliedern. Vorsitzender ist Stefan Klebert, Finanzvorstand
seit dem 1.~November 2025 Alexander Kocherscheidt, der Bernd Brinker abl\"oste; Johannes
Giloth ist Chief Operating Officer, Dr.~Nadine Sterley verantwortet People \&
Sustainability. Die drei \"ubrigen Mitglieder f\"uhren je eine Division: Peter Lauwers
Pharma \& Food Applications, Kai Becker Pure Flow Processing, Klaus Stojentin Nutrition
Plant Engineering\Q{GeaGB}{264}. Alle drei Divisionsposten wurden zum 1.~Januar 2026 neu
besetzt; sie sind die personelle Seite der im zweiten Halbjahr 2025 vorgestellten neuen
Konzernstruktur\Q{GeaGB}{12}. Vorsitzender des Aufsichtsrats ist Prof.~Dieter Kempf,
stellvertretender Vorsitzender Rainer Gr\"obel\Q{GeaGB}{265}.

\bk{Dass Finanzvorstand und alle drei Divisionsvorst\"ande innerhalb von zwei Monaten
wechseln, ist der gr\"o{\ss}te Umbau der Organe seit dem Amtsantritt des Vorsitzenden. Der
Vorsitzende selbst bleibt, und mit ihm die Strategie; die Ergebnisse der n\"achsten zwei
Jahre entstehen aber unter weitgehend neuer Leitung. F\"ur die Beurteilung der
Prognosetreue in Abschnitt~\ref{sec:strategie} hei{\ss}t das: die belegte Bilanz geh\"ort
der alten Mannschaft.}

\subsection{Eigent\"umer und Gesch\"afte der Organe}

Zum 31.12.2025 bestanden \Mio{\AktienMio}{Aktien} Stimmrechtsaktien gegen\"uber
\Mio{\AktienVorjahr}{Aktien} ein Jahr zuvor; der R\"uckgang um \AktienRueckgang{} stammt aus dem
Einzug zur\"uckgekaufter Aktien\Q{GeaGB}{14}. Die Aktion\"arsanalyse vom Dezember 2025
identifizierte \Pz{\IdentifiziertAnteil} aller Aktien: \Pz{\Institutionelle} institutionelle
Anleger, \Pz{\Privatanleger} Privatanleger, \Pz{\Grossaktionaere} Gro{\ss}aktion\"are. Der
Streubesitz betr\"agt \Pz{\Streubesitz}\Q{GeaGB}{14}. Einen kontrollierenden Aktion\"ar gibt
es nicht.

In den zw\"olf Monaten bis zum Berichtsstand meldete GEA \DdKaeufe{} K\"aufe von
\DdPersonen{} Vorstandsmitgliedern zu Kursen zwischen \je{\DdPreisVon}{EUR} und
\je{\DdPreisBis}{EUR}; ein Verkauf wurde nicht gefunden. Der Vorstandsvorsitzende kaufte in
zwei Tranchen f\"ur zusammen \je{\DdKlebertVolumen}{EUR}\QW{EQS News, Directors' Dealings der
GEA Group}{quelle:ddgea}{19.09.2026}.

\bk{Sechs K\"aufe, kein Verkauf, vier Personen, alle innerhalb weniger Wochen nach einem
Kursr\"uckgang: das ist ein ungew\"ohnlich geschlossenes Bild. Es belegt, dass die Organe den
Kurs von Mai und Juni 2026 f\"ur zu niedrig hielten. Diese K\"aufe lagen zwischen
\je{\DdPreisVon}{EUR} und \je{\DdPreisBis}{EUR}; der heutige Kurs liegt dar\"uber. Ein
Organkauf ist ein Signal \"uber einen Preis, nicht \"uber jeden Preis.}

\subsection{Sicht des Kapitalmarkts}

Moody's best\"atigte im Juni 2025 das Rating Baa1 mit stabilem Ausblick und hob die
konservative Finanzierungspolitik und die geringe Verschuldung hervor\Q{GeaGB}{16}. Die
Dividendenpolitik sieht seit dem Kapitalmarkttag 2024 eine Aussch\"uttungsquote von rund
\Pz{\AusschuettungsZiel} des Konzernergebnisses vor; f\"ur 2025 schlagen Vorstand und
Aufsichtsrat \je{\Dividende}{EUR} je Aktie vor nach \je{\DividendeVorjahr}{EUR}\Q{GeaGB}{16}.
Das entspricht \Mio{\DividendeSummeKuenftig}{EUR} und bei einem Konzernergebnis von
\Mio{\Jahresergebnis}{EUR} einer Quote von \AusschuettungsQuote{}; die Politik ist damit
eingehalten.

\subsection{Analystenkonsens}

Der Bericht selbst weist die Verteilung der Empfehlungen zum 16.~Februar 2026 aus:
\Pz{\AnalystPositiv} positiv, \Pz{\AnalystNeutral} neutral, \Pz{\AnalystNegativ}
negativ\Q{GeaGB}{15}. Diese Angabe ist zum Berichtsstand sieben Monate alt und wird deshalb
durch eine Sekund\"arquelle erg\"anzt: \KonsensAnalysten{} Analysten, Gesamtvotum
\enquote{Outperform}, mittleres Kursziel \je{\KonsensZiel}{EUR}, h\"ochstes
\je{\KonsensHoch}{EUR}, niedrigstes \je{\KonsensTief}{EUR} bei einem Schlusskurs von
\je{\KonsensKurs}{EUR}\QW{MarketScreener, Konsens der Analysten zur GEA
Group}{quelle:konsens}{19.09.2026}. Das mittlere Kursziel liegt \Pz{\KonsensAufschlag} \"uber
dem Kurs.

\bk{Der Konsens ist kein Beleg f\"ur einen Wert, sondern eine Tatsache \"uber die Erwartung
anderer. Bemerkenswert ist weniger sein Niveau als seine Streuung: zwischen
\je{\KonsensTief}{EUR} und \je{\KonsensHoch}{EUR} liegt fast die H\"alfte des Kurses. Eine so
weite Spanne bei einem stetigen Gesch\"aft zeigt, dass der Streit nicht dem laufenden Jahr
gilt, sondern der Frage, wie lange das Wachstum anh\"alt. Genau diese Frage
\"ubersetzt Abschnitt~\ref{sec:votum} in eine Zahl.}

\subsection{Nicht verf\"ugbare Angaben}

Zum Berichtsstand nicht beschafft: die Einzelkursziele der \KonsensAnalysten{} Analysten (nur
Mittelwert und Spanne liegen vor), die Leerverkaufsquote, und Marktanteile je Division. GEA
nennt keine Marktanteile; die Aussage der Wettbewerbsposition bleibt deshalb ohne diese
Zahl (Abschnitt~\ref{sec:moat}). Die Vollst\"andigkeit der Directors'-Dealings-Meldungen ist
nicht garantiert: die Meldungen wurden einzeln abgerufen, nicht aus einem amtlichen Register
exportiert.

\subsection{Wettbewerbsposition: ein Burggraben ist eine Zahl oder ein Adjektiv}
\label{sec:moat}

\begin{table}[htbp]\centering\small
\caption{Behauptete Vorteile und die Zahl, die sie zeigen w\"urde.}\label{tab:moat}
\begin{tabularx}{\linewidth}{@{}XXl@{}}\toprule
Behauptung (Quelle) & Zahl, die sie zeigen w\"urde & H\"alt?\\\midrule
\enquote{leading market positions as a mechanical and plant engineering company for the
energy-intensive food, beverage and pharmaceutical industries}\Q{GeaGB}{24} &
Marktanteil je Division & nicht belegbar, GEA nennt keinen\\
Vernetzte installierte Basis als Grundlage des Servicegesch\"afts\Q{GeaGB}{24} &
Serviceanteil am Umsatz 2025 & ja: \Pz{\Serviceanteil}\Q{GeaGB}{2}\\
Preissetzungsmacht & EBITDA-Marge im Zeitverlauf & ja: Tabelle~\ref{tab:reihe}\\
Kapitaleffizienz & ROCE 2025 & ja: \Pz{\Roce}\Q{GeaGB}{2}\\
\bottomrule\end{tabularx}\end{table}

\bk{Von vier Behauptungen sind drei mit einer Zahl belegt und eine nicht. Der belastbarste
Beleg ist der Serviceanteil: \Pz{\Serviceanteil} des Umsatzes stammen aus dem Gesch\"aft an
Maschinen, die bereits beim Kunden stehen. Das ist wiederkehrender Umsatz mit h\"oherer
Marge und geringerem Investitionsbedarf, und es erkl\"art, warum die Marge seit 2020
in jedem einzelnen Jahr gestiegen ist. Es belegt keinen Marktanteil; wer den wissen will,
findet ihn im Bericht nicht.}
